% !TeX root = ../../../main.tex

Das \texttt{analysis}-Modul enthält die Funktionen, die die Logdateien mit pandas auswerten.
Die Funktion \texttt{start\_analysis()} wird durch die \texttt{main} aufgerufen startet die einzelnen Analyseschritte.
Dafür wird die Logdatei eingelesen und in einem DataFrame gespeichert.

Durch die Funktion \texttt{\_print()} werden die Ergebnisse in die Konsole geschrieben.
Die Funktion in Codeausschnitt~\ref{lst:sortNumeric} implementiert eine numerische Sortierung.
Hierzu werden die ersten beiden Ebenen des MultiIndex betrachtet.
Sind die Indexwerte Strings, wird mithilfe eines regulären Ausdrucks der darin enthaltene numerische Wert
ausgelesen und in einen \texttt{float} umgewandelt.
Das \texttt{if hasattr} dient lediglich als Absicherung, dass nur in Strings die Zahl extrahiert wird, während bei
numerischen Werten direkt sortiert wird.
Dadurch werden die Ergebnisse in der numerisch korrekten Reihenfolge und nicht lexikografisch sortiert.
Wenn der Einfluss eines bestimmten Parameters betrachtet wird, steht dieser in dem ersten Level des Index, sodass die
Sortierung zunächst anhand dieses Parameters und anschließend anhand des zweiten Levels erfolgt.

\begin{lstlisting}[style=pythoncode, caption={Numerische Sortierung der Ergebnisse}, label={lst:sortNumeric}]
def _sort_numeric(data):
    return data.sort_index(
        level=[0, 1],
        key=lambda x: x.str.extract(r"([\d.]+)")[0].astype(float)
		if hasattr(x, 'str') else x
    )
\end{lstlisting}

Für die Auswertung stehen vier Funktionen zur Verfügung, die den Einfluss der Parameter Concurrency, Delay, Retry Count
und Strategie untersuchen.
Zusätzlich existiert eine Funktion, welche die Ergebnisse unabhängig von einem dieser Parameter auswertet.
Den vier parameterabhängigen Analysefunktionen wird jeweils der Parameter übergeben, dessen Einfluss auf die
Versuchsergebnisse untersucht werden soll.
Dieser Parameter bildet bei der anschließenden Auswertung mittels \texttt{groupby} die erste Ebene des Index und
besitzt somit die höchste Priorität bei der Sortierung und Auswertung der Ergebnisse.
Die Analysefunktionen rufen anschließend die einzelnen Funktionen zur Berechnung der in Abschnitt~\ref{sec:purpose}
beschriebenen Kenngrößen auf.
Dadurch kann dieselbe Berechnungslogik für unterschiedliche Parameter verwendet werden, ohne die eigentliche Berechnung
der Kenngrößen mehrfach implementieren zu müssen.

Im Folgendem wird nur der Fall, ohne zusätzliche Parameter betrachtet, da sich lediglich das Auslesen und die
Gruppierung nach diesem ändert.
Des Weiteren enthält jede Funktion, außer \texttt{\_throughput} den Parameter \texttt{pstatus}, welcher entweder auf
\texttt{success}, \texttt{failed} oder mit \texttt{""} leer gesetzt werden kann.
Dadurch kann das DataFrame wahlweise mit allen Einträgen, ausschließlich mit erfolgreichen Einträgen oder
ausschließlich mit fehlgeschlagenen Einträgen ausgegeben werden.

Die \texttt{\_rates()}-Funktion berechnet die Erfolgs- und Fehlerrate der Experimente in Prozent.
Hierfür wird das DataFrame zunächst nach dem Status gruppiert und die Anzahl der Einträge je Status gezählt und
gespeichert.
Die jeweilige Rate ergibt sich anschließend aus der Division von der Anzahl der Einträge und der Gesamtanzahl der
Einträge.
Das wird am Schluss mit 100 multipliziert, damit ein Prozentwert ausgegeben wird und nicht die Dezimalzahl.

Für die \texttt{\_retry\_distribution} wird das DataFrame nach \texttt{retries} und dem Status gruppiert.
Durch \texttt{dropna=False} werden \texttt{None}- beziehungsweise \texttt{NaN}-Werte berücksichtigt.
Dabei zählt \texttt{size()} die Einträge und \texttt{unstack()} wandelt die zweite Indexebene in Spalten um, während
\texttt{fill\_value=0} dafür sorgt, dass nicht vorhandene Kombinationen mit 0 dargestellt werden.

Bei der Funktion \texttt{\_execution\_comparison} werden Ausführungszeiten verglichen.
Dafür wird das DataFrame nach dem Status gruppiert und für jede Gruppe der Durchschnitt mit \texttt{mean()} von der
\texttt{total\_ms} berechnet.
Die Quartile in \texttt{\_quartiles} berechnen sich analog dazu, außer dass \texttt{quantile([0.50, 0.95, 0.99, 0.999])}
anstelle von \texttt{mean} verwendet wurde.

Zur Bestimmung des durch den Retry-Mechanismus entstehenden Overheads werden zunächst die einzelnen Ausführungszeiten
der Transaktionsversuche betrachtet.
Es wird aus \texttt{attempts\_ms} jeweils der erste Eintrag entfernt, da dieser die ursprüngliche Ausführung der
Transaktion und somit keinen Retry darstellt.
Aus den verbleibenden Ausführungszeiten wird anschließend mit \texttt{min()} die kürzeste Ausführungszeit bestimmt und
in \texttt{best\_ms} gespeichert.
Der Overhead ergibt sich schließlich aus der Differenz zwischen der gesamten Ausführungszeit \texttt{total\_ms} und der
kürzesten gemessenen Ausführungszeit.
Anschließend wird aus dem \texttt{overhead} der Durchschnitt berechnet.
Dadurch wird der zusätzliche Zeitaufwand bestimmt, der über die reine Ausführungszeit des schnellsten
Transaktionsversuchs hinausgeht.

Für die Berechnung des Durchsatzes wird zunächst aus der \texttt{run\_id} die jeweilige Experiment-ID extrahiert und in
der Spalte \texttt{e} gespeichert.
Anschließend werden die Einträge anhand dieser Experiment-ID gruppiert.
Für jedes Experiment werden dabei der früheste Startzeitpunkt und der späteste Endzeitpunkt aller Transaktionen
bestimmt.
Die Differenz zwischen diesen beiden Zeitpunkten ergibt die gesamte Laufzeit des jeweiligen Experiments.
Zusätzlich werden ausschließlich die erfolgreich abgeschlossenen Transaktionen betrachtet und pro Experiment gezählt.
Die Anzahl der erfolgreichen Transaktionen wird anschließend mit der zuvor bestimmten Laufzeit ins Verhältnis gesetzt.
Der resultierende Wert \texttt{throughput} gibt somit an, wie viele erfolgreiche Transaktionen pro Zeiteinheit
innerhalb des jeweiligen Experiments abgeschlossen wurden.

Durch diese Funktionen lassen sich die relevanten Kenngrößen berechnen und strukturiert nach den jeweiligen
Parametern darstellen.
Dadurch können die Ergebnisse gezielt miteinander verglichen und Rückschlüsse auf den Einfluss der
einzelnen Parameter gezogen werden.
